\documentclass[11pt]{article}
\usepackage[margin=1in]{geometry}
\usepackage{amsmath,amssymb,amsthm,graphicx,microtype}
\usepackage[hidelinks]{hyperref}
\usepackage[T1]{fontenc}

\newtheorem{theorem}{Theorem}
\newcommand{\N}{\mathbb N}
\newcommand{\ellTwo}{\ell_2(\N)}

\title{An Alternating Coordinate-Block Counterexample to a Banach
Feichtinger Conjecture}
\author{Discovery packet for arXiv:2301.03365}
\date{13 August 2026}

\begin{document}
\maketitle

\begin{center}
\fbox{\parbox{0.92\linewidth}{\textbf{Status.}
Full counterexample, pending human review. Conjecture~2.5 is false under the
strict diagonal-dominance formulation in Definition~2.4 and also under the
usual Riesz-sequence reading. The construction is not localized and leaves
the paper's Theorem~2.7 intact.}}
\end{center}

\section{Source conjecture and terminology}

Krishna's \emph{Localized Bounded Below Approximate Schauder Frames are
Finite Unions of Approximate Riesz Sequences}, arXiv:2301.03365
\cite{Krishna}, calls a pair $(\{f_n\},\{\tau_n\})$ an approximate Schauder
frame (ASF) when
\[
 Sx=\sum_{n=1}^{\infty} f_n(x)\tau_n
\]
defines a bounded invertible operator. It is \emph{bounded below} when
$\inf_n|f_n(\tau_n)|>0$.

Definition~2.4 asks for a finite partition
$\N=Q_1\cup\cdots\cup Q_r$ such that every class obeys
\begin{equation}\label{eq:dd}
 \inf_{n\in Q_j}\left(
 |f_n(\tau_n)|-\sum_{\substack{m\in Q_j\\m\ne n}}
 |f_n(\tau_m)|\right)>0.
\end{equation}
Conjecture~2.5 states that every bounded-below ASF admits the proposed finite
decomposition. The conjecture itself prints ``ARBs,'' which is not defined in
the paper, while the title, abstract, section heading, Definition~2.4, and
Theorem~2.7 use ``ARSs.'' The proof of Theorem~2.7 establishes exactly
\eqref{eq:dd}. We therefore refute this precise displayed condition and also
check the ordinary Riesz-sequence interpretation.

\begin{figure}[ht]
\centering
\includegraphics[width=0.8\linewidth]{figures/open_problem_crop.png}
\caption{Definition~2.4 and Conjecture~2.5 on page 2 of the source paper.}
\label{fig:source}
\end{figure}

\section{Proof intuition}

An ASF only requires the \emph{ordered} rank-one series to converge to an
invertible operator. This permits large finite cancellation patterns whose
supports escape to infinity. On the $k$th coordinate, use an odd block of
$2k+1$ copies of the same vector, with coefficient functionals carrying
alternating signs. The block sum is one coordinate projection, so the full
frame operator is the identity. But the block sizes are unbounded. Any finite
coloring repeats a color inside some block; that same-color pair has
off-diagonal interaction of modulus one, equal to the diagonal, and destroys
strict diagonal dominance. The identical vectors also rule out a Riesz
sequence independently of the source's terminology.

\section{The counterexample}

Let $(e_k)_{k\ge1}$ be the standard orthonormal basis of $\ellTwo$, with
coordinate functionals $e_k^*$. Partition $\N$ into the consecutive odd
blocks
\[
 I_k=\{k^2,k^2+1,\ldots,(k+1)^2-1\},\qquad |I_k|=2k+1.
\]
For $n\in I_k$, define
\begin{equation}\label{eq:construction}
 \tau_n=e_k,
 \qquad
 f_n=(-1)^{n-k^2}e_k^*.
\end{equation}

\begin{theorem}\label{thm:main}
The pair in \eqref{eq:construction} is a bounded-below ASF for $\ellTwo$ with
frame operator the identity, but no finite partition of $\N$ satisfies
\eqref{eq:dd}. Moreover, no finite partition makes every vector subsequence
$(\tau_n)_{n\in Q_j}$ a Riesz sequence.
\end{theorem}

\begin{proof}
Write $P_k$ for the projection onto
$\operatorname{span}\{e_1,\ldots,e_k\}$ and $P_0=0$. For $x\in\ellTwo$ and
$n\in I_k$,
\[
 f_n(x)\tau_n=(-1)^{n-k^2}x_ke_k.
\]
The signs on $I_k$ are $+1,-1,+1,\ldots,-1,+1$. Consequently their total is
one, and every initial sum within that block is either zero or one. It follows
that every frame partial sum ending inside $I_k$ equals either $P_{k-1}x$ or
$P_kx$. Since $P_kx\to x$ in $\ell_2$, the ordered series converges for every
$x$ and
\[
 Sx=\sum_{n=1}^{\infty}f_n(x)\tau_n=x.
\]
Thus $S=I$ is bounded and invertible. Also
\[
 |f_n(\tau_n)|=1\qquad(n\in\N),
\]
so the ASF is bounded below with lower value one.

Now consider an arbitrary partition $\N=Q_1\cup\cdots\cup Q_r$. Choose $k$
with $2k+1>r$. By the pigeonhole principle there exist distinct
$n,m\in I_k$ in one class, say $Q_j$. Since $\tau_n=\tau_m=e_k$,
\[
 |f_n(\tau_n)|=1,
 \qquad |f_n(\tau_m)|=1.
\]
Therefore
\[
 |f_n(\tau_n)|-
 \sum_{\substack{q\in Q_j\\q\ne n}}|f_n(\tau_q)|\le0,
\]
and the infimum required in \eqref{eq:dd} cannot be positive. Since the
partition was arbitrary, no finite partition works.

Finally, the same indices have identical vectors $\tau_n=\tau_m$. The
synthesis map on their class sends the nonzero finitely supported coefficient
vector with entries $1$ at $n$ and $-1$ at $m$ to zero. Hence that vector
subsequence is not a Riesz sequence (or Riesz basic sequence) under the usual
injectivity/minimality requirement. No finite partition can avoid this
collision in all blocks.
\end{proof}

\section{Scope and adversarial checks}

The example does not satisfy intrinsic localization. Indeed, within $I_k$ one
has $|f_n(\tau_m)|=1$ for all $n,m$, while the block diameters $2k$ tend to
infinity. Thus no fixed $A>0$ and $s>1$ can give the paper's decay bound
$A/(1+|n-m|)^s$. This is why the counterexample refutes the general
Conjecture~2.5 without contradicting its localized Theorem~2.7.

The construction also exposes the exact role of conditional ordering. The
partial sums are uniformly bounded coordinate projections, so the ASF
definition is met without ambiguity, but alternating signs allow arbitrarily
large finite redundancy. Stronger hypotheses such as unconditionality or
factorization through an $\ell_p$ analysis space are not claimed here.

The script \texttt{verify\_finite\_blocks.py} checks the block sums and partial
sums exactly and exhausts small coloring cases. It is a sanity check; the
general proof is the argument above.

\section{Novelty and review recommendation}

A bounded search on 13 August 2026 covered the local run indexes and parsed
arXiv corpus; exact conjecture and title phrases; the ASF/ARS terminology;
and web searches for Conjecture~2.5, arXiv:2301.03365 with
``counterexample/refute/false,'' and bounded-below ASF counterexamples. The
results found the target and mirrors but no prior counterexample or explicit
resolution. This does not establish exhaustive novelty.

A human reviewer should focus on (i) whether ordered convergence is indeed
the intended ASF convention in Definition~2.1; (ii) the scope chosen in light
of the source's `ARB' typo; and (iii) the bounded novelty search. Human review
has not yet been completed.

\begin{thebibliography}{9}
\bibitem{Krishna}
K. Mahesh Krishna,
\emph{Localized Bounded Below Approximate Schauder Frames are Finite Unions
of Approximate Riesz Sequences}, arXiv:2301.03365 (2023).
\end{thebibliography}

\end{document}
